% from ega4-13-fr.tex
\oldpage[IV-3]{191}
\begin{definition}[13.2.2]
\label{IV.13.2.2-fr}
Sous les hypothèses de (13.2.1), on dit que $f$ est équidimensionnel au point
$x$ (ou que $X$ est équidimensionnel sur $Y$ au point $x$) si
\[
  \dim_x(f^{-1}(f(x)))=\dim(f^{-1}(\eta)).
\]
On dit que $f$ est équidimensionnel (ou que $X$ est équidimensionnel sur $Y$)
si $f$ est équidimensionnel en tout point $x\in X$.
\end{definition}

Il résulte du théorème de Chevalley (13.1.3) que l'ensemble des $x\in X$ où
$f$ est équidimensionnel est \emph{ouvert non vide}. En outre, si $f$ est
équidimensionnel au point $x$, toutes les composantes irréductibles de
$f^{-1}(f(x))$ qui contiennent $x$ ont \emph{même dimension}, car chacune
d'elles a une dimension qui est $\geq\dim(f^{-1}(\eta))$ (13.1.6) et
$\leq\dim_x(f^{-1}(f(x)))$ en vertu de ($\mathbf 0$, 14.1.5).

\begin{proposition}[13.2.3]
\label{IV.13.2.3-fr}
Soient $Y$ un préschéma irréductible localement noethérien, $X$ un préschéma
irréductible, $f:X\to Y$ un morphisme dominant localement de type fini~;
soient $\eta$ le point générique de $Y$, $x$ un point de $X$, $y=f(x)$, et
supposons que l'on ait
\[
\tag{13.2.3.1}
\label{IV.13.2.3.1-fr}
  \dim(\mathcal O_x)=\dim(\mathcal O_y)
  +\dim(\mathcal O_x\otimes_{\mathcal O_y}k(y)).
\]
Alors $f$ est équidimensionnel au point $x$. La réciproque est vraie si les
deux membres de l'inégalité (5.6.5.2) sont égaux, en particulier si
$\mathcal O_y$ est universellement caténaire.
\end{proposition}

\begin{proof}
Cela résulte aussitôt de (5.6.5.2) et de l'inégalité $\delta(x)\geq0$
(13.1.1).
\end{proof}

\begin{env}[13.2.4]
\label{IV.13.2.4-fr}
Soient maintenant, de façon générale, $Y$ un préschéma localement noethérien,
$f:X\to Y$ un morphisme de type fini, $x$ un point de $X$, $X'$ une partie
fermée irréductible de $X$ contenant $x$, $y=f(x)$,
$Y'=\overline{f(X')}$, $\eta'$ le point générique de $Y'$. Désignons encore
par $X'$ et $Y'$ les sous-préschémas fermés réduits de $X$ et $Y$
respectivement ayant $X'$ et $Y'$ pour espaces sous-jacents~; la restriction
$X'\to Y$ de $f$ se factorise donc en
$X'\xrightarrow{f'}Y'\xrightarrow{j}Y$, où $j$ est l'injection canonique
(\textbf{I}, 5.2.2), et $f'$ est de type fini (1.5.4). Posons alors pour
abréger
\[
\tag{13.2.4.1}
\label{IV.13.2.4.1-fr}
  A=\mathcal O_{Y,y},\quad B=\mathcal O_{X,x},\quad
  A'=\mathcal O_{Y',y},\quad B'=\mathcal O_{X',x}.
\]
La formule (5.6.5.2) appliquée au morphisme dominant $f'$ et aux préschémas
irréductibles $X'$, $Y'$, donne
\[
\tag{13.2.4.2}
\label{IV.13.2.4.2-fr}
  \dim(B')\leq\dim(A')+\dim(B'\otimes_{A'}k(y))
  -(\dim_x(f'^{-1}(y))-\dim(f'^{-1}(\eta'))).
\]
D'autre part, l'anneau local $A'$ (resp. $B'$, $B'\otimes_{A'}k(y)$) est un
anneau quotient de $A$ (resp. $B$, $B\otimes_Ak(y)$), donc
($\mathbf 0$, 16.1.2.1) on a
\[
\tag{13.2.4.3}
\label{IV.13.2.4.3-fr}
  \dim(A')\leq\dim(A),\quad \dim(B')\leq\dim(B),\quad
  \dim(B'\otimes_{A'}k(y))\leq\dim(B\otimes_Ak(y)).
\]
On déduit donc en premier lieu de (13.2.4.3) et ($\mathbf 0$, 16.3.9)
\[
\tag{13.2.4.4}
\label{IV.13.2.4.4-fr}
  \dim(B')\leq\dim(A')+\dim(B'\otimes_{A'}k(y))
  \leq\dim(A)+\dim(B\otimes_Ak(y)).
\]
En outre, en vertu de ($\mathbf 0$, 16.3.9), on a aussi les inégalités
\[
\tag{13.2.4.5}
\label{IV.13.2.4.5-fr}
  \dim(B')\leq\dim(B)\leq\dim(A)+\dim(B\otimes_Ak(y)).
\]
\end{env}

